\chapter{Structs}\label{ch:structs}
\index{struct}

Every interesting program eventually needs to group related data together.
An array of coordinates, a pair of first-and-last names, a record that bundles
a user's email with their role---the moment your data grows beyond a single
value, you reach for a \emph{struct}.

Lattice structs are lightweight, named containers for fields. They are
pass-by-value, they support methods through trait implementations and callable
fields, and they integrate naturally with the phase system you met in
\cref{ch:phases-explained}. In this chapter we will define structs, read and
write their fields, attach behavior, customise equality, and peer under the
hood with the reflection API.

%% ─────────────────────────────────────────────────────────────────────
\section{Defining and Instantiating Structs}\label{sec:struct-defining}
\index{struct!definition}

A struct definition gives a name and a list of typed fields:

\begin{lstlisting}[language=Lattice, caption={A minimal struct definition}]
struct Point { x: any, y: any }
\end{lstlisting}

The keyword \lstinline|struct| is followed by a name that must begin with an
uppercase letter, then a pair of braces enclosing comma-separated fields.
Each field has a name and a type annotation. Right now the most common
annotation is \lstinline|any|, which accepts every type; we will explore
richer type annotations when we reach generics in \cref{ch:generics}.

\begin{defnbox}{Struct}
A \emph{struct} is a named, composite value type whose data is stored in
a fixed set of named fields. Structs are \textbf{pass-by-value}: assigning
a struct to a new variable produces an independent copy.
\end{defnbox}

Once defined, you create an instance with a \emph{struct literal}---the struct
name followed by braces and field initialisers:

\begin{lstlisting}[language=Lattice, caption={Creating struct instances}]
struct Point { x: any, y: any }

let origin = Point { x: 0, y: 0 }
let target = Point { x: 42, y: 17 }

print(origin)  // Point { x: 0, y: 0 }
print(target)  // Point { x: 42, y: 17 }
\end{lstlisting}

Every field must be supplied when constructing an instance---there are no
default values at the language level. If you want defaults, write a factory
function:

\begin{lstlisting}[language=Lattice, caption={A factory function with defaults}]
struct Config {
    host: String,
    port: Int,
    verbose: Bool
}

fn default_config() -> Config {
    return Config {
        host: "localhost",
        port: 8080,
        verbose: false
    }
}

let cfg = default_config()
print(cfg.port)  // 8080
\end{lstlisting}

\subsection{Nested Structs}

Fields can hold any value, including other structs:

\begin{lstlisting}[language=Lattice, caption={Nested structs}]
struct Point { x: any, y: any }
struct Rect { origin: any, width: any, height: any }

let r = Rect {
    origin: Point { x: 5, y: 10 },
    width: 100,
    height: 50
}

print(r.origin.x)  // 5
print(r.width)     // 100
\end{lstlisting}

Nesting is a natural way to compose data. Because structs are
pass-by-value, the \lstinline|Point| inside the \lstinline|Rect| is a
full copy, not a reference. Modifying the original \lstinline|Point| after
constructing the \lstinline|Rect| has no effect on \lstinline|r.origin|.

\subsection{Structs with Typed Fields}

Although \lstinline|any| works everywhere, you can annotate fields with
concrete types for documentation and, in strict mode, static checking:

\begin{lstlisting}[language=Lattice, caption={Typed struct fields}]
struct Entity {
    id: Int,
    name: String,
    x: Int,
    y: Int,
    hp: Int,
    max_hp: Int,
    tag: String
}
\end{lstlisting}

This reads like a blueprint: an \lstinline|Entity| carries an integer ID, a
name, a position, health, and a tag string. Readers of your code---and the
compiler in strict mode---will thank you for the specificity.

%% ─────────────────────────────────────────────────────────────────────
\section{Field Access and Mutation}\label{sec:struct-fields}
\index{struct!field access}
\index{struct!mutation}

Reading a field uses dot notation:

\begin{lstlisting}[language=Lattice, caption={Field access}]
struct Point { x: any, y: any }

let p = Point { x: 3, y: 4 }
print(p.x)      // 3
print(p.y)      // 4
print(p.x + p.y) // 7
\end{lstlisting}

\subsection{Mutation Requires \texttt{flux}}

Because structs follow the phase system, a struct bound with \lstinline|let|
(or \lstinline|fix|) is \emph{crystallised}---its fields cannot be changed.
To mutate fields, bind the struct with \lstinline|flux|:

\begin{lstlisting}[language=Lattice, caption={Mutating struct fields}]
struct Point { x: any, y: any }

flux p = Point { x: 1, y: 2 }
print(p.x)  // 1

p.x = 99
print(p.x)  // 99

p.y = 88
print(p.y)  // 88
\end{lstlisting}

Compound assignment operators work on fields too:

\begin{lstlisting}[language=Lattice, caption={Compound field assignment}]
struct Point { x: any, y: any }

flux p = Point { x: 10, y: 20 }
p.x += 5
print(p.x)  // 15
p.y -= 10
print(p.y)  // 10
p.x *= 2
print(p.x)  // 30
\end{lstlisting}

\begin{warnbox}{Mutation is Shallow}
When you mutate a field on a \lstinline|flux| struct, you are updating
\emph{that copy's} field. Because structs are pass-by-value, any other
variable that received a copy of the struct is unaffected. This is one of
the most important mental-model points in Lattice. We will return to it in
\cref{sec:struct-pass-by-value}.
\end{warnbox}

\subsection{Optional Field Access}

If you have a value that might be \lstinline|nil|, the optional-chaining
operator \lstinline|?.| lets you access fields safely:

\begin{lstlisting}[language=Lattice, caption={Optional field access}]
struct User { name: any, email: any }

fn find_user(id: any) -> any {
    if id == 1 {
        return User { name: "Alice", email: "alice@example.com" }
    }
    return nil
}

let user = find_user(42)
print(user?.name)   // nil (no crash)
print(user?.email)  // nil
\end{lstlisting}

If the receiver is \lstinline|nil|, the entire expression evaluates to
\lstinline|nil| instead of producing a runtime error. This pairs well
with the \lstinline|??| nil-coalescing operator covered in
\cref{ch:control-flow}.

\subsection{Passing Structs to Functions}

Structs can be passed as function arguments and returned from functions.
Remember: each pass produces an independent copy.

\begin{lstlisting}[language=Lattice, caption={Structs as arguments and return values}]
struct Point { x: any, y: any }

fn sum_point(pt: any) -> any {
    return pt.x + pt.y
}

fn make_point(a: any, b: any) -> Point {
    return Point { x: a, y: b }
}

let p = Point { x: 3, y: 4 }
print(sum_point(p))  // 7

let q = make_point(100, 200)
print(q.x)  // 100
print(q.y)  // 200
\end{lstlisting}

\subsection{Structs in Collections}

Arrays, maps, and sets can all hold structs:

\begin{lstlisting}[language=Lattice, caption={An array of structs}]
struct Point { x: any, y: any }

let points = [
    Point { x: 1, y: 2 },
    Point { x: 3, y: 4 },
    Point { x: 5, y: 6 }
]

for pt in points {
    print(pt.x + pt.y)
}
// Output: 3, 7, 11
print(len(points))  // 3
\end{lstlisting}

%% ─────────────────────────────────────────────────────────────────────
\section{Methods via Traits: \texttt{impl} Blocks}\label{sec:struct-methods}
\index{struct!methods}
\index{impl}

Lattice attaches behaviour to structs through \emph{trait implementations}.
A trait declares a set of method signatures; an \lstinline|impl| block
provides the bodies for a specific struct. We will explore traits in depth
in \cref{ch:traits}, but here is enough to get you started.

\begin{lstlisting}[language=Lattice, caption={A trait and its implementation}]
struct Point { x: any, y: any }

trait Describable {
    fn describe(self: any) -> any
}

impl Describable for Point {
    fn describe(self: any) -> any {
        return "Point(" + to_string(self.x) + ", " + to_string(self.y) + ")"
    }
}

let p = Point { x: 7, y: 8 }
print(p.describe())  // Point(7, 8)
\end{lstlisting}

The first parameter of every method is \lstinline|self|, which receives
a copy of the struct instance the method was called on. Inside the body
you read fields with \lstinline|self.x|, \lstinline|self.y|, and so on.

\begin{notebox}{Methods Return New Values}
Because \lstinline|self| is a copy, methods that ``modify'' the struct
must return a new instance. The original is untouched:

\begin{lstlisting}[language=Lattice]
trait Scalable {
    fn scale(self: any, factor: any) -> any
}

impl Scalable for Point {
    fn scale(self: any, factor: any) -> any {
        return Point { x: self.x * factor, y: self.y * factor }
    }
}

let p = Point { x: 5, y: 10 }
let q = p.scale(3)
print(q.x)  // 15
print(q.y)  // 30
print(p.x)  // 5  (unchanged)
\end{lstlisting}
\end{notebox}

\subsection{Multiple Traits per Struct}

A single struct can implement as many traits as you like:

\begin{lstlisting}[language=Lattice, caption={A struct implementing multiple traits}]
struct Point { x: any, y: any }

trait HasArea {
    fn area(self: any) -> any
}

trait Describable {
    fn describe(self: any) -> any
}

impl HasArea for Point {
    fn area(self: any) -> any {
        return self.x * self.y
    }
}

impl Describable for Point {
    fn describe(self: any) -> any {
        return "Point(" + to_string(self.x) + ", " + to_string(self.y) + ")"
    }
}

let p = Point { x: 3, y: 4 }
print(p.area())      // 12
print(p.describe())  // Point(3, 4)
\end{lstlisting}

Under the hood, the compiler registers each method with a key of the form
\texttt{TypeName::method\_name} in the global scope. When you call
\lstinline|p.describe()|, the VM looks up \texttt{Point::describe} and
invokes it with \lstinline|self| bound to a clone of \lstinline|p|. You can
see how method dispatch is compiled in \filename{src/stackcompiler.c}---look
for the \texttt{OP\_INVOKE} family of opcodes.

\subsection{Methods in Loops}

Methods compose naturally with iteration:

\begin{lstlisting}[language=Lattice, caption={Calling methods in a loop}]
struct Point { x: any, y: any }

trait HasArea {
    fn area(self: any) -> any
}

impl HasArea for Point {
    fn area(self: any) -> any {
        return self.x * self.y
    }
}

let points = [
    Point { x: 1, y: 1 },
    Point { x: 2, y: 3 },
    Point { x: 4, y: 5 }
]

flux total = 0
for p in points {
    total = total + p.area()
}
print(total)  // 1 + 6 + 20 = 27
\end{lstlisting}

%% ─────────────────────────────────────────────────────────────────────
\section{Callable Struct Fields}\label{sec:struct-callable-fields}
\index{struct!callable fields}
\index{closure!as struct field}

Besides attaching methods through traits, you can embed closures directly
as struct fields. Any field whose value is a closure becomes \emph{callable}:
when you invoke it with dot syntax, Lattice finds the closure field and
calls it.

\begin{lstlisting}[language=Lattice, caption={A struct with a callable closure field}]
struct Greeter {
    greeting: String,
    greet: Fn
}

let hello = Greeter {
    greeting: "Hello",
    greet: |self| {
        self.greeting + ", world!"
    }
}

print(hello.greet())  // Hello, world!
\end{lstlisting}

The closure receives \lstinline|self| as its first argument when called
through dot notation, just like a trait method. The difference is that the
closure is stored \emph{inside} the struct value itself---it travels with
the data.

\subsection{Stateful Callable Fields}

Callable fields shine when you want per-instance behaviour. Consider a
vending machine where each operation is a closure:

\begin{lstlisting}[language=Lattice, caption={Vending machine with callable fields}]
struct VendingMachine {
    balance: Int,
    inventory: Map,
    state: String,
    insert_coin: Fn,
    select_item: Fn,
    dispense: Fn,
    display: Fn
}

fn make_vending_machine() -> VendingMachine {
    flux inv = Map::new()
    inv.set("cola", 150)
    inv.set("chips", 100)
    inv.set("candy", 75)

    return VendingMachine {
        balance: 0,
        inventory: inv,
        state: "idle",
        insert_coin: |self, amount| {
            self.balance + amount
        },
        select_item: |self, item| {
            if !self.inventory.has(item) {
                "error: unknown item"
            } else {
                let price = self.inventory.get(item)
                if self.balance < price {
                    "error: insufficient funds"
                } else {
                    "ok:" + item
                }
            }
        },
        dispense: |self, item| {
            let price = self.inventory.get(item)
            let change = self.balance - price
            "Dispensing " + item + "! Change: " + to_string(change) + " cents"
        },
        display: |self| {
            "Balance: " + to_string(self.balance) + " cents"
        }
    }
}
\end{lstlisting}

Each closure can read \lstinline|self.balance| and
\lstinline|self.inventory|. The caller updates the struct's mutable state
after each operation:

\begin{lstlisting}[language=Lattice, caption={Driving the vending machine}]
flux vm = make_vending_machine()
print(vm.display())  // Balance: 0 cents

vm.balance = vm.insert_coin(25)
vm.balance = vm.insert_coin(25)
vm.balance = vm.insert_coin(25)
print(vm.display())  // Balance: 75 cents

let result = vm.select_item("candy")
print(result)  // ok:candy
print(vm.dispense("candy"))  // Dispensing candy! Change: 0 cents
\end{lstlisting}

\begin{tipbox}{Trait Methods vs.\ Callable Fields}
Use \textbf{trait methods} when the behaviour is shared across multiple
structs that implement the same interface. Use \textbf{callable fields}
when the behaviour is unique to each \emph{instance}, or when you want to
swap behaviour at runtime by assigning a different closure to the field.
\end{tipbox}

\subsection{How the Runtime Resolves Calls}

When you write \lstinline|vm.display()|, the runtime checks two places:

\begin{enumerate}
  \item \textbf{Callable field}: Does the struct have a field named
    \lstinline|display| whose value is a closure? If so, call it, passing
    the struct as \lstinline|self|.
  \item \textbf{Trait method}: Is there an impl method registered as
    \lstinline|VendingMachine::display|? If so, call that with a clone
    of the struct as \lstinline|self|.
\end{enumerate}

Callable fields take priority. If a struct has both a closure field named
\lstinline|area| and a trait method \lstinline|area|, the closure field
wins. You can inspect how this dispatch works in \filename{src/eval.c}---the
method-call handler first scans \lstinline|field_names| for a matching
closure before falling through to the impl registry.

%% ─────────────────────────────────────────────────────────────────────
\section{Custom Equality with \texttt{eq}}\label{sec:struct-equality}
\index{struct!equality}
\index{eq (custom equality)}

By default, two structs are equal if they have the same name and all their
fields are recursively equal:

\begin{lstlisting}[language=Lattice, caption={Default struct equality}]
struct Point { x: any, y: any }

let a = Point { x: 1, y: 2 }
let b = Point { x: 1, y: 2 }
let c = Point { x: 3, y: 4 }

print(a == b)  // true
print(a == c)  // false
print(a != c)  // true
\end{lstlisting}

This is structural equality---it compares data, not identity. Two distinct
instances with the same field values are considered equal.

\subsection{Overriding Equality with an \texttt{eq} Field}

Sometimes you want a different notion of equality. Perhaps two users should
be equal if they share the same ID, regardless of other fields. Lattice
lets you override equality by adding a closure field named \lstinline|eq|:

\begin{lstlisting}[language=Lattice, caption={Custom equality via the \texttt{eq} field}]
struct User {
    id: Int,
    name: String,
    email: String,
    eq: Fn
}

fn make_user(id: Int, name: String, email: String) -> User {
    return User {
        id: id,
        name: name,
        email: email,
        eq: |self, other| {
            self.id == other.id
        }
    }
}

let alice1 = make_user(1, "Alice", "alice@example.com")
let alice2 = make_user(1, "Alice A.", "alice.a@example.com")
let bob = make_user(2, "Bob", "bob@example.com")

print(alice1 == alice2)  // true  (same ID)
print(alice1 == bob)     // false (different ID)
\end{lstlisting}

The \lstinline|eq| closure receives two arguments: \lstinline|self| (the
left-hand operand) and \lstinline|other| (the right-hand operand). Its
return value is converted to a boolean. The \lstinline|!=| operator
automatically negates the result.

\begin{warnbox}{The \texttt{eq} Field Must Be a Closure}
The custom equality mechanism specifically checks for a field named
\lstinline|eq| whose value is a \lstinline|VAL_CLOSURE|. If you name a
non-closure field \lstinline|eq|, it will not be used for equality
comparisons---the default structural equality applies instead.
\end{warnbox}

\subsection{When Custom Equality Matters}

Custom equality is valuable for:

\begin{itemize}
  \item \textbf{Identity-based comparison}: Compare by ID or key rather
    than every field.
  \item \textbf{Ignoring computed fields}: Skip over cached or derived
    data that should not affect equality.
  \item \textbf{Tolerance-based comparison}: For structs containing
    floating-point values, you might compare with a tolerance:
\end{itemize}

\begin{lstlisting}[language=Lattice, caption={Approximate equality for geometry}]
struct Coordinate {
    lat: Float,
    lng: Float,
    eq: Fn
}

fn make_coord(lat: Float, lng: Float) -> Coordinate {
    return Coordinate {
        lat: lat,
        lng: lng,
        eq: |self, other| {
            let dlat = self.lat - other.lat
            let dlng = self.lng - other.lng
            if dlat < 0 { dlat = 0.0 - dlat }
            if dlng < 0 { dlng = 0.0 - dlng }
            dlat < 0.0001 && dlng < 0.0001
        }
    }
}

let a = make_coord(51.5074, -0.1278)
let b = make_coord(51.5074, -0.1278)
print(a == b)  // true
\end{lstlisting}

%% ─────────────────────────────────────────────────────────────────────
\section{Struct Reflection}\label{sec:struct-reflection}
\index{struct!reflection}
\index{struct\_name}
\index{struct\_fields}
\index{struct\_to\_map}
\index{struct\_from\_map}

Lattice provides four built-in functions for inspecting and transforming
structs at runtime. These live in the category of \emph{reflection}---code
that examines the shape of data programmatically.

\subsection{\texttt{struct\_name()}}

Returns the type name of a struct as a string:

\begin{lstlisting}[language=Lattice, caption={Getting a struct's type name}]
struct Point { x: any, y: any }
struct Rect { width: any, height: any }

let p = Point { x: 10, y: 20 }
let r = Rect { width: 5, height: 3 }

print(struct_name(p))  // "Point"
print(struct_name(r))  // "Rect"
\end{lstlisting}

This is useful for logging, serialisation, and dynamic dispatch.

\subsection{\texttt{struct\_fields()}}

Returns an array of field-name strings:

\begin{lstlisting}[language=Lattice, caption={Listing a struct's fields}]
struct Config {
    host: String,
    port: Int,
    verbose: Bool
}

let cfg = Config { host: "localhost", port: 8080, verbose: false }
let fields = struct_fields(cfg)
print(fields)  // ["host", "port", "verbose"]
\end{lstlisting}

The order of names matches the order in the struct definition.

\subsection{\texttt{struct\_to\_map()}}

Converts a struct into a \lstinline|Map| where keys are field names and
values are field values:

\begin{lstlisting}[language=Lattice, caption={Converting a struct to a map}]
struct Point { x: any, y: any }

let p = Point { x: 42, y: 17 }
let m = struct_to_map(p)

print(m)            // {x: 42, y: 17}
print(m.get("x"))   // 42
print(m.get("y"))   // 17
\end{lstlisting}

This is the fastest way to bridge between structs and the map-based world
of JSON serialisation (see \cref{ch:json-toml-yaml-csv}).

\subsection{\texttt{struct\_from\_map()}}

Goes the other direction---creates a struct from a type name and a map:

\begin{lstlisting}[language=Lattice, caption={Creating a struct from a map}]
struct Point { x: any, y: any }

flux m = Map::new()
m.set("x", 100)
m.set("y", 200)

let p = struct_from_map("Point", m)
print(p.x)  // 100
print(p.y)  // 200
\end{lstlisting}

If a field is missing from the map, it defaults to \lstinline|nil|. The
struct definition must already exist in scope---\lstinline|struct_from_map|
looks up the registered struct metadata at runtime.

\subsection{A Practical Reflection Example}

Combining these functions, you can write a generic ``debug print'' utility:

\begin{lstlisting}[language=Lattice, caption={A generic struct inspector}]
struct Point { x: any, y: any }
struct User { name: any, age: any }

fn inspect(val: any) -> String {
    let name = struct_name(val)
    let fields = struct_fields(val)
    let m = struct_to_map(val)
    flux parts = []
    for f in fields {
        parts.push(f + "=" + to_string(m.get(f)))
    }
    return name + "(" + parts.join(", ") + ")"
}

let p = Point { x: 3, y: 4 }
let u = User { name: "Alice", age: 30 }

print(inspect(p))  // Point(x=3, y=4)
print(inspect(u))  // User(name=Alice, age=30)
\end{lstlisting}

\begin{notebox}{Reflection at Runtime}
Struct reflection functions operate at \emph{runtime}. The compiler stores
field metadata as a hidden global (keyed as \texttt{\_\_struct\_StructName}
in the constant pool). This means reflection incurs a small cost, but it
works for any struct, even those defined in imported modules. See
\filename{src/stackcompiler.c} for how metadata is emitted during
compilation.
\end{notebox}

%% ─────────────────────────────────────────────────────────────────────
\section{Struct Destructuring}\label{sec:struct-destructuring}
\index{struct!destructuring}

When you want to extract multiple fields at once, destructuring is more
concise than repeated dot access:

\begin{lstlisting}[language=Lattice, caption={Struct destructuring}]
struct Point { x: any, y: any }

let p = Point { x: 55, y: 66 }
let { x, y } = p
print(x)  // 55
print(y)  // 66
\end{lstlisting}

The \lstinline|let { x, y } = p| syntax binds local variables
\lstinline|x| and \lstinline|y| to the values of
\lstinline|p.x| and \lstinline|p.y| respectively. The variable names must
match the field names.

Destructuring works in any binding context where patterns are accepted,
including \lstinline|match| arms (see \cref{ch:pattern-matching}).

\begin{tipbox}{Selective Destructuring}
You do not have to destructure every field. Extracting only the fields you
need is perfectly valid:

\begin{lstlisting}[language=Lattice]
struct Rect { origin: any, width: any, height: any }

let r = Rect { origin: Point { x: 0, y: 0 }, width: 100, height: 50 }
let { width, height } = r
print(width * height)  // 5000
\end{lstlisting}
\end{tipbox}

%% ─────────────────────────────────────────────────────────────────────
\section{Pass-by-Value Semantics}\label{sec:struct-pass-by-value}
\index{struct!pass-by-value}
\index{pass-by-value}

This is the section that will save you from the most common struct-related
confusion. Lattice structs are \textbf{pass-by-value}. Every assignment,
every function argument, every return value produces an \emph{independent
copy} of the struct. There are no hidden references, no shared state, no
aliasing surprises.

\begin{lstlisting}[language=Lattice, caption={Pass-by-value in action}]
struct Point { x: any, y: any }

flux a = Point { x: 1, y: 2 }
flux b = a          // b is a copy of a
b.x = 999

print(a.x)  // 1   (unchanged)
print(b.x)  // 999
\end{lstlisting}

After \lstinline|flux b = a|, the variables \lstinline|a| and \lstinline|b|
hold entirely separate structs. Mutating one has zero effect on the other.

\subsection{Implications for Function Calls}

When you pass a struct to a function, the function receives its own copy:

\begin{lstlisting}[language=Lattice, caption={Functions receive copies}]
struct Counter { value: any }

fn try_to_increment(c: any) {
    flux local = c
    local.value = local.value + 1
    print(local.value)  // 11
}

let counter = Counter { value: 10 }
try_to_increment(counter)
print(counter.value)  // 10 (unchanged)
\end{lstlisting}

The function's local copy is incremented, but the caller's
\lstinline|counter| retains its original value. If you want the change to
be visible to the caller, the function must \emph{return} the new struct:

\begin{lstlisting}[language=Lattice, caption={Returning the modified struct}]
struct Counter { value: any }

fn increment(c: any) -> Counter {
    return Counter { value: c.value + 1 }
}

flux counter = Counter { value: 10 }
counter = increment(counter)
print(counter.value)  // 11
\end{lstlisting}

\subsection{Why Pass-by-Value?}

You might wonder why Lattice chose value semantics when many popular
languages use references. The answer ties into the phase system:

\begin{itemize}
  \item \textbf{No aliasing bugs}: Two variables never point at the same
    struct, so freezing one cannot affect another.
  \item \textbf{Thread safety}: In concurrent code (\cref{ch:scope-spawn}),
    spawned tasks receive their own copy of any struct, eliminating data
    races by construction.
  \item \textbf{Predictability}: When you see \lstinline|let x = y|, you
    know that \lstinline|x| is independent. No need to check whether the
    type is a ``value type'' or a ``reference type.''
\end{itemize}

\begin{defnbox}{Value Semantics}
In Lattice, structs (and buffers) are \textbf{value types}. Assignment copies
the entire struct, including all nested data. The runtime uses
\texttt{value\_deep\_clone()} in \filename{include/value.h} to produce the
copy. For large structs this means more memory, but it guarantees isolation.
\end{defnbox}

\subsection{The \texttt{clone} Keyword}

If you want to be explicit about copying---perhaps for readability in
performance-sensitive code---you can use the \lstinline|clone| keyword:

\begin{lstlisting}[language=Lattice, caption={Explicit cloning}]
struct Point { x: any, y: any }

let original = Point { x: 1, y: 2 }
flux copy = clone original
copy.x = 42

print(original.x)  // 1
print(copy.x)      // 42
\end{lstlisting}

In practice, assignment already clones, so \lstinline|clone| is primarily
a documentation signal: ``I know this is making a copy, and I intend that.''

\subsection{When Copies Get Expensive}

For small structs like \lstinline|Point|, copying is negligible. For structs
with many fields or large nested arrays, the cost can add up. Strategies for
managing this include:

\begin{itemize}
  \item \textbf{Keep structs lean}: Store IDs or indices into a central
    collection rather than embedding large data directly.
  \item \textbf{Use \lstinline|Ref| for shared state}: When you genuinely
    need shared mutable access, wrap a value in a \lstinline|Ref|
    (covered in \cref{ch:phases-explained}).
  \item \textbf{Prefer functional updates}: Return new structs from methods
    rather than mutating fields---the compiler can often optimise this path.
\end{itemize}

%% ─────────────────────────────────────────────────────────────────────
\section{A Larger Example: Entity Component System}\label{sec:struct-ecs}
\index{ECS}

Let's bring everything together with a complete example. An Entity Component
System (ECS) is a game architecture pattern where entities are bags of
data and systems are functions that transform arrays of entities. It is a
perfect showcase for struct-centric design.

\begin{lstlisting}[language=Lattice, caption={An ECS in Lattice}]
struct Entity {
    id: Int,
    name: String,
    x: Int,
    y: Int,
    dx: Int,
    dy: Int,
    hp: Int,
    max_hp: Int,
    tag: String
}

fn make_entity(id: Int, name: String, tag: String) -> Entity {
    return Entity {
        id: id, name: name,
        x: 0, y: 0, dx: 0, dy: 0,
        hp: 0, max_hp: 0, tag: tag
    }
}

fn render_entity(e: Entity) {
    flux line = "  [" + to_string(e.id) + "] " + e.name
    line += " at (" + to_string(e.x) + "," + to_string(e.y) + ")"
    if e.max_hp > 0 {
        line += " HP:" + to_string(e.hp) + "/" + to_string(e.max_hp)
    }
    print(line)
}

// System: apply velocity to position
fn movement_system(entities: [Entity]) -> [Entity] {
    flux result = []
    for e in entities {
        flux moved = e
        moved.x = e.x + e.dx
        moved.y = e.y + e.dy
        result.push(moved)
    }
    return result
}
\end{lstlisting}

The movement system iterates over entities, creates updated copies with new
positions, and returns the new array. Because structs are pass-by-value,
each \lstinline|flux moved = e| produces an independent copy that we can
mutate without affecting the original.

\begin{lstlisting}[language=Lattice, caption={Running the ECS simulation}]
fn main() {
    flux player = make_entity(0, "Hero", "player")
    player.x = 0
    player.y = 0
    player.dx = 1
    player.dy = 1
    player.hp = 100
    player.max_hp = 100

    flux goblin = make_entity(1, "Goblin", "enemy")
    goblin.x = 3
    goblin.y = 3
    goblin.dx = -1
    goblin.dy = -1

    flux entities = [player, goblin]

    flux tick = 0
    while tick < 3 {
        print("--- Tick " + to_string(tick) + " ---")
        for e in entities {
            render_entity(e)
        }
        entities = movement_system(entities)
        tick += 1
    }
}
\end{lstlisting}

This pattern---immutable-by-default data flowing through pure transformation
functions---is idiomatic Lattice. The phase system and value semantics make
it safe; the struct system makes it readable.

%% ─────────────────────────────────────────────────────────────────────
\section{Exercises}\label{sec:struct-exercises}

\begin{enumerate}
  \item \textbf{RGB colour struct.}
    Define a struct \lstinline|Color| with fields \lstinline|r|,
    \lstinline|g|, and \lstinline|b| (all integers 0--255). Write a
    function \lstinline|mix(a: Color, b: Color) -> Color| that returns
    the average of each channel. Verify that mixing
    \lstinline|Color { r: 255, g: 0, b: 0 }| with
    \lstinline|Color { r: 0, g: 0, b: 255 }| gives
    \lstinline|Color { r: 127, g: 0, b: 127 }|.

  \item \textbf{Bank account.}
    Create a \lstinline|BankAccount| struct with fields \lstinline|owner|,
    \lstinline|balance|, and callable fields \lstinline|deposit| and
    \lstinline|withdraw|. The \lstinline|withdraw| field should return
    an error string if the balance is insufficient. Drive a sequence of
    transactions and print the final balance.

  \item \textbf{Custom equality.}
    Define a \lstinline|Student| struct with fields \lstinline|id|,
    \lstinline|name|, and \lstinline|gpa|. Add an \lstinline|eq| closure
    that compares students by \lstinline|id| only. Demonstrate that two
    \lstinline|Student| instances with the same ID but different GPAs are
    considered equal.

  \item \textbf{Reflection round-trip.}
    Write a function that takes any struct, converts it to a map with
    \lstinline|struct_to_map|, doubles every numeric value in the map,
    and converts it back with \lstinline|struct_from_map|. Test it with
    \lstinline|Point { x: 5, y: 10 }| and verify the result is
    \lstinline|Point { x: 10, y: 20 }|.

  \item \textbf{Value semantics drill.}
    Without running the code, predict the output of this program, then
    verify your prediction:
    \begin{lstlisting}[language=Lattice]
struct Box { value: any }
flux a = Box { value: [1, 2, 3] }
flux b = a
b.value.push(4)
print(len(a.value))
print(len(b.value))
    \end{lstlisting}
\end{enumerate}

%% ─────────────────────────────────────────────────────────────────────
\section*{What's Next}

Structs give your data a name and a shape. But what if a value can be
\emph{one of several shapes}? A network response might be a success or
a failure. A drawing command might be a circle, a rectangle, or a line.
In the next chapter we meet \emph{enums}---Lattice's answer to
sum types---and discover how to model these ``one of many'' scenarios
with precision and safety.
